\section{Predictive Results for Architecture Choice}

\paragraph{Reading guide.}
Each result below has the same structure. First, the topology fixes what each worker or coordinator can
observe. Second, that visibility pattern induces a cost, error, or planning bound through a simple
inequality such as a union bound, the data processing inequality, or a critical-path argument. The
epistemic notation is bookkeeping for visibility: if an agent cannot observe a fact directly, recovering
that fact later requires communication, compression, or review.

\paragraph{Worked examples.}
Each theorem is paired with a small illustrative agent-architecture example. These examples are design
calculations, not benchmark measurements.

\subsection{Reviewer Gates and Correlated Failure}

A minimal reviewer gate can be modeled as a coherent feed-forward loop: a generator proposes an action,
but the action is released only if a reviewer also emits approval.

\begin{theorem}[Reviewer-Gated Error Suppression]
\label{thm:reviewer-gate}
Let $A_{\mathrm{gen}}$ be a generator agent and $A_{\mathrm{ver}}$ be a reviewer. Let
$E_{\mathrm{gen}}$ denote the event that the generator emits an erroneous candidate, and let
$M_{\mathrm{ver}}$ denote the event that the reviewer misses that error and still emits approval.
\begin{enumerate}[leftmargin=*]
\item[(a)] \textbf{Direct link.} Without a reviewer gate,
\[
P(\mathrm{Fail}_{\mathrm{direct}}) = P(E_{\mathrm{gen}}).
\]
\item[(b)] \textbf{Reviewer gate under independence.} If $E_{\mathrm{gen}}$ and $M_{\mathrm{ver}}$ are independent,
\[
P(\mathrm{Fail}_{\mathrm{gate}}) = P(E_{\mathrm{gen}})\,P(M_{\mathrm{ver}}).
\]
\item[(c)] \textbf{Reviewer gate under correlation.} Let $p = P(E_{\mathrm{gen}})$, $q = P(M_{\mathrm{ver}})$,
and let $\rho$ be the Bernoulli correlation between the two events. Then
\[
P(\mathrm{Fail}_{\mathrm{gate}})
= pq + \rho\sqrt{p(1-p)q(1-q)}.
\]
\end{enumerate}
\end{theorem}

\paragraph{Assumptions.}
The theorem uses the minimal failure model for a two-stage execution gate: in the direct case an
erroneous action is emitted exactly when the generator emits an erroneous candidate; in the gated case an
erroneous action is emitted exactly when the generator emits an error and the reviewer still approves it.

\paragraph{Operational takeaway.}
Adding a reviewer buys multiplicative safety only when the reviewer is not making the same mistakes as
the generator. The gate provides the conjunction; diversity in model family, evidence source, or checking
mechanism determines how much of that multiplicative gain survives in practice.

\begin{proof}[Proof sketch]
In the direct topology there is only one path from the generator to the action sink, so
$\mathrm{Fail}_{\mathrm{direct}} = E_{\mathrm{gen}}$. In the gated topology an erroneous final action can
occur only if the generator emits an error and the reviewer misses it while still approving, so
$\mathrm{Fail}_{\mathrm{gate}} = E_{\mathrm{gen}} \wedge M_{\mathrm{ver}}$. Independence gives the product
$pq$. For the correlated case, writing $X = \mathbf{1}_{E_{\mathrm{gen}}}$ and
$Y = \mathbf{1}_{M_{\mathrm{ver}}}$ yields
\[
\mathbb{E}[XY] = \mathbb{E}[X]\mathbb{E}[Y] + \mathrm{Cov}(X,Y)
= pq + \rho\sqrt{p(1-p)q(1-q)}.
\]
\end{proof}

\paragraph{Worked Agent Example (Execution Gate).}
Suppose a code-writing agent proposes a migration and a reviewer checks it before execution. If the
generator error rate is $p=0.1$ and the reviewer miss rate is $q=0.1$, independence gives
$P(\mathrm{Fail}_{\mathrm{gate}})=0.01$ instead of $0.1$. But if reviewer misses are positively correlated
with generator errors, say $\rho=0.5$, the failure rate rises to $0.055$. The gate still helps, but much
less than the independent ideal.

\subsection{Error Amplification in Independent and Centralized Aggregation}

\begin{theorem}[Error Amplification Bound]
\label{thm:error-amplification}
Consider $n$ agents processing independent subtasks with error probability $p$, aggregated by a final
combiner. The error amplification factor $A$ (aggregate failure probability divided by single-agent
failure probability) depends on the coordination topology:
\begin{enumerate}[leftmargin=*]
\item[(a)] \textbf{Independent workers.} If the aggregator observes only final worker outputs,
\begin{equation}
A_{\otimes} \le n.
\label{eq:error-amp-independent}
\end{equation}
\item[(b)] \textbf{Centralized coordination.} If a hub sees all worker outputs before aggregation and detects
an error with rate $d = P(\text{hub detects error} \mid \text{error occurred})$, then
\begin{equation}
A_{\circ} \le n(1-d).
\label{eq:error-amp-centralized}
\end{equation}
\end{enumerate}
\end{theorem}

\paragraph{Assumptions.}
Worker errors are independent across subtasks, the aggregate fails when at least one erroneous subtask
survives to the final output, and in the centralized case the hub has an approximately stationary
per-error detection rate $d$ across subtasks.

\paragraph{Operational takeaway.}
Adding workers without an effective review bottleneck scales failure opportunities roughly with the number
of workers. A hub helps exactly to the extent that it can see and suppress cross-worker inconsistencies.

\begin{proof}[Proof sketch]
For independent workers, the aggregate fails with probability
$P(\mathrm{failure}_{\otimes}) = 1 - (1-p)^n \le np$ by the union bound, so
$A_{\otimes} = P(\mathrm{failure}_{\otimes})/p \le n$. With a centralized hub, each subtask error reaches
the final output only if it occurs and escapes detection, so
$P(\mathrm{failure}_{\circ}) = 1 - (1-p(1-d))^n \le np(1-d)$, giving
$A_{\circ} \le n(1-d)$. In the small-$p$ regime the exact ratio approaches $1/(1-d)$, recovering the
intuition that review quality directly controls the gain from centralization.
\end{proof}

\paragraph{Worked Agent Example (Code Review Gate).}
Suppose three code-generation workers each draft part of a deployment change, and a release gate accepts
the merged result. If each worker has error rate $p=0.1$, an ungated independent aggregate fails with
probability $1-(1-0.1)^3 = 0.271$, so $A_{\otimes}=2.71$. If a central reviewer catches half of all
worker errors ($d=0.5$), then the centralized failure rate becomes
$1-(1-0.1(1-0.5))^3 \approx 0.143$, so $A_{\circ}\approx 1.43$.

\subsection{Sequential Coordination Penalty}

\begin{theorem}[Sequential Coordination Penalty]
\label{thm:sequential-penalty}
For a task with $k$ strictly ordered steps decomposed across agents, each inter-agent handoff must at
minimum establish $K_{\mathrm{receiver}}(\mathrm{result}_j)$. The overhead is:
\begin{enumerate}[leftmargin=*]
\item[(a)] \textbf{Single agent.} $K_i(\mathrm{result}_j)$ is trivial because the agent computed the result.
Total cost is $k \cdot c_{\mathrm{step}}$.
\item[(b)] \textbf{Multi-agent pipeline.} Each of $h \le k-1$ handoffs incurs communication cost
$c_{\mathrm{comm}}$ and epistemic reconstruction loss
\begin{equation}
\Delta I_j = I(S_{\mathrm{sender}}; r_j) - I(\mathrm{obs}_{\mathrm{receiver}}(S_{\mathrm{sender}}); r_j),
\label{eq:epistemic-reconstruction-loss}
\end{equation}
which is nonnegative by data processing and strictly positive whenever the handoff is lossy on the
task-relevant support.
\end{enumerate}
\end{theorem}

\paragraph{Assumptions.}
The task is strictly sequential, later steps depend on earlier results, each handoff transmits only a
summary of the sender's relevant state, and the receiver must reconstruct enough of that state to
continue execution.

\paragraph{Operational takeaway.}
This is the warning against gratuitous decomposition: a sequential pipeline benefits from multiple agents
only if a handoff creates new observations or new capabilities. Otherwise the system pays communication
cost plus lossy-summary cost.

\begin{proof}[Proof sketch]
A single agent keeps prior results in its own local state, so no handoff is required and total cost is
$k \cdot c_{\mathrm{step}}$. When step $j$ is handed from agent $a$ to agent $b$, the receiver observes
only a transformed summary of the sender's state. By the data processing inequality,
$I(\mathrm{obs}_b(S_a); r_j) \le I(S_a; r_j)$, so the receiver cannot gain information at handoff and
typically loses some. Each handoff therefore adds message-construction cost, parsing cost, and
reconstruction effort proportional to the lost task-relevant information.
\end{proof}

\paragraph{Worked Agent Example (Bug-Fix Relay).}
Consider a three-step bug-fix task: interpret the failing test, locate the root cause, and write the
patch. A single coding agent pays $3c_{\mathrm{step}}$. If the task is split across planner, debugger,
and patcher with $h=2$ handoffs, $c_{\mathrm{step}}=100$ tokens, $c_{\mathrm{comm}}=25$, and
$c_{\mathrm{recon}}=20$ per handoff, then $C_{\mathrm{single}}=300$ while
$C_{\mathrm{multi}}=3\cdot100+2\cdot(25+20)=390$.

\subsection{Parallel Acceleration Under Task Decomposability}

\begin{theorem}[Parallel Acceleration under Epistemic Independence]
\label{thm:parallel-acceleration}
A task decomposes into $m$ subtasks. Define two subtasks as \emph{epistemically independent} iff
neither solution requires knowledge of the other's result.
\begin{enumerate}[leftmargin=*]
\item[(a)] Under full epistemic independence with centralized coordination, the speedup is
\begin{equation}
S = \frac{\sum_{i=1}^{m} c_{\mathrm{sub}_i}}{\max_i(c_{\mathrm{sub}_i}) + c_{\mathrm{assign}} + c_{\mathrm{agg}}}.
\label{eq:parallel-speedup}
\end{equation}
\item[(b)] For propositions determined solely by the tuple of worker outputs, the coordinator attains the
corresponding pooled-output knowledge at aggregation as an architectural byproduct. Error detection from
Theorem~\ref{thm:error-amplification} therefore comes with little extra communication cost.
\item[(c)] If subtasks are only partially independent, parallel execution sacrifices quality in proportion
to the unresolved mutual information between subtask solutions.
\end{enumerate}
\end{theorem}

\paragraph{Assumptions.}
Subtasks can be assigned independently, coordinator assignment and aggregation are additive overhead
terms, and no hidden blocking dependency forces a worker to wait for another worker's intermediate
result.

\paragraph{Operational takeaway.}
This is the positive case for multi-agent systems: if subtasks can be solved from local observations
alone, specialization plus a coordinator can reduce wall-clock cost and often improve quality. If
subtasks share unresolved dependencies, parallelism turns into premature decomposition.

\begin{proof}[Proof sketch]
Under epistemic independence, all $m$ subtasks can run concurrently and wall-clock cost is determined by
the slowest subtask plus coordinator overhead, giving Eq.~\eqref{eq:parallel-speedup}. Because the
coordinator must receive all worker outputs to aggregate them, it also obtains pooled-output visibility as
part of normal execution. When subtasks are not independent, some worker decisions depend on facts hidden
in another worker's local state; running them in parallel then replaces needed information sharing with
guesswork, reducing solution quality.
\end{proof}

\paragraph{Worked Agent Example (Due-Diligence Swarm).}
Suppose an orchestrator assigns three largely independent subtasks for a vendor assessment: legal review,
security posture, and cost modeling. If each subtask costs about $8$ minutes of agent time and
assignment plus aggregation costs $2$ minutes total, then
\[
S = \frac{8+8+8}{8+2} = 2.4.
\]
The coordinator also receives all three outputs at aggregation, so cross-checking across workstreams
comes with little extra communication cost.

\subsection{Tool Density and the Coordination Tax}

\begin{theorem}[Tool Density Scaling]
\label{thm:tool-density}
Consider $t$ tools distributed across $n$ agents, with $|T_i| = t/n$ tools per agent under an
approximately balanced partition.
\begin{enumerate}[leftmargin=*]
\item[(a)] \textbf{Single agent.} Planning cost is $O(t)$ per step because the agent scans one unified tool
inventory in context.
\item[(b)] \textbf{Multi-agent.} If agent $i$ needs a tool in $T_j$ for $j \ne i$, it must discover the
tool, request remote execution, and interpret the result with information loss $\Delta I$ from
Eq.~\eqref{eq:epistemic-reconstruction-loss}. The coordination overhead per step scales as
\begin{equation}
C_{\mathrm{coord}} =
O\!\left(t \cdot \left(1 - \frac{|T_i \cap T_{\mathrm{needed}}|}{|T_{\mathrm{needed}}|}\right) \cdot c_{\mathrm{comm}}\right).
\label{eq:tool-coord-overhead}
\end{equation}
\item[(c)] \textbf{Cross-agent planning overhead.} The planner must also reason about remote tool
capabilities, yielding worst-case planning cost
\begin{equation}
C_{\mathrm{plan}} = O(tn),
\label{eq:tool-planning-cost}
\end{equation}
which is a multiplicative $n$-factor over the single-agent baseline and becomes bilinear when both $t$
and $n$ grow.
\end{enumerate}
\end{theorem}

\paragraph{Assumptions.}
Tools are partitioned approximately evenly across agents, remote tool use requires explicit discovery and
delegation, and the planner may in the worst case need to reason over every tool-agent pairing relevant
to the current subgoal.

\paragraph{Operational takeaway.}
Splitting a toolset across agents changes one local decision into three coordination steps: discover who
has the tool, delegate execution, then interpret the returned result without the executor's full context.
That is why tool-heavy tasks often punish over-distribution.

\begin{proof}[Proof sketch]
A single agent scans one tool inventory, so planning cost is linear in $t$. Under partition, a worker
sees only its local tool subset. Remote-tool use therefore introduces discovery, delegation, and result
reconstruction costs, producing the coordination term in Eq.~\eqref{eq:tool-coord-overhead}. Planning
also becomes second-order: the agent must reason not just about tool applicability but about which other
agent knows how to apply which tool. In the worst case that produces $O(tn)$ candidate comparisons.
\end{proof}

\paragraph{Worked Agent Example (Tool-Split SWE Agent).}
Imagine a software-engineering assistant with $18$ tools split across $3$ agents: one owns search tools,
one owns git and test tools, and one owns deployment tools. A planner diagnosing a failing CI job
directly holds only $6$ tools and must treat the remaining $12$ as remote. The planning problem expands
from ``which of 18 tools should I call next?'' to ``which agent owns the relevant tool, how do I route
the subproblem, and how do I interpret the result without that agent's full execution context?''

% Source map:
% - 5.1 adapted from article/sections/04-operad.tex theorem on CFFL.
% - 5.2--5.5 adapted from article/sections/06-discussion.tex predictive theorems.
